%قبل از شروع مسئله ترجیحاً چیزی ننویسید. اما اگر خواستید بنویسید، در قالب لیستی (فقرات و شمارش) بنویسید.
%برای اضافه کردن نمره، بعدی از شروع‌مسئله، [x] را اضافه کنید و به جای x  نمره مورد نظر را قرار دهید. برای زیرمسئله‌ها هم به همین صورت است.
\شروع‌مسئله[18]
برای هر یک از قسمت‌های این سوال، مسئله را به کمک گراف مدل کرده و سپس به کمک الگوریتم مناسب حل کنید. جزئیات مدل و دلیل درستی الگوریتم خود را شرح دهید.
%اگر مسئله چند قسمت داشت به این صورت عمل کنید:
\شروع{زیرمسائل}
\زیرمسئله[6]
در یک کشور بین برخی از شهر‌ها جاده وجود دارد به طوری که از هر شهر می‌توان به شهر‌های دیگر رفت. این جاده‌ها می‌توانند یک طرفه یا دو طرفه باشند. همچنین هنگام ورود به هر شهر باید  مالیاتی  مانند $r > 0$    پرداخت شود. در این مسئله یک مسافر می‌خواهد از شهر $A$  به شهر $B$ برود و    کمترین هزینه را پرداخت کند. برای یافتن مسیر بهینه به او کمک کنید.
 % مسئلۀ ۱آ

\پاسخ{
	برای هر شهر یک رأس در نظر می‌گیریم.
	اگر از شهر 
	$A$
	به شهر 
	$B$
	جاده‌ای وجود داشته باشد، از رأس متناظر $A$ به راس متناظر $B$ یک یال جهت‌دار رسم می‌کنیم. اگر هم بین دو شهر جاده دوطرفه باشد، یک یال جهت‌دار از 
 $A$ به  $B$	
و یک یال جهت‌دار هم از 
 $B$ به  $A$
رسم می‌کنیم. حال به یک مسئلۀ جست‌و‌جو در گراف رسیده‌ایم. اما در این مسئله رأس‌ها وزن دارند و نه یال‌ها. وزن‌های رأس‌ها را نیز قدرمطلق مالیات آن شهر می‌گیریم. برای این‌که وزن‌ها را از رأس‌ها به یال‌ها منتقل کنیم، کافی است به تمام یال‌های ورودی هر شهر، وزن آن شهر را نسبت دهیم. پس به یک مسئله ساده جست‌و‌جو کوتاه‌ترین  مسیر در گراف وزن‌دار (با وزن نامنفی) رسیده‌ایم که به کمک الگوریتم دایسترا قابل حل است.

}



\زیرمسئله[12]
یک شیمی‌دان می‌تواند در هر فرآیند یک ماده را به ماده‌ای دیگر تبدیل کند. به صورت دقیق‌تر برای برخی از مواد این شیمی‌دان می‌داند که می‌توان یک گرم از ماده اول را به 
$r$
گرم ($r > 0$) از ماده دوم (در یک مرحله) تبدیل کند. او می‌خواهد بداند که مستقل از تعداد آزمایش‌ها، یک گرم مس را حداکثر به چند گرم طلا می‌توان تبدیل کرد.



\پاسخ{
	برای هر ماده یک رأس در نظر می‌گیریم. همچنین اگر یک گرم از مادهٔ
	$A$
	در یک آزمایش به  
	$r$ گرم از مادهٔ
	$B$
	تبدیل شود، آنگاه یک یال جهت‌دار از 
	$A$
	به 
	$B$
	با وزن
	$- \log r$
	رسم می‌کنیم. ادعا می‌کنیم مسئله به یافتن کوتاه‌ترین مسیر از مس به طلا در این گراف وزن‌دار تبدیل شده است. برای این موضوع کافی است توجه کنیم که کوتاه‌ترین مسیر در حقیقت
	$- \sum_i \log r_i$
	را بین این دو رأس کمینه می‌کند و به عبارت دیگر 
$\sum_i \log r_i = \log(\prod_i r_i)$
را  بیشینه می‌کند. پس مسیری را خروجی می‌دهد که 
$\prod_i r_i$
را بیشینه می‌کند که برابر با مقدار طلایی هست که از یک گرم مس در یک مسیر حاصل می‌شود که بیشینه کردن این مقدار برابر با پیدا کردن مسیری است که بیشترین طلا را از یک گرم مس حاصل می‌کند. اگر $w$ وزن کوتاه‌ترین مسیر باشد، مقدار بیشترین طلای حاصل شده $2^{-w}$ است.


برای حل این مسئله گراف می‌توانیم از الگوریتم 
بلمن-فورد
استفاده کنیم. توجه کنیم که یال‌ها در این مسئله می‌توانند وزن منفی داشته باشند. حال اگر مسیری از مس به طلا نباشد، نمی‌توان چنین تبدیلی انجام داد.
همچنین اگر بتوان با دور منفی به این طلا رسید، مقدار طلایی که می‌توان تبدیل کرد نامحدود خواهد بود و در غیر این صورت هم که می‌توانیم
$2^{-w}$
را خروجی دهیم. 
(البته توجه کنیم که اگر قانون بقای جرم را در حل مسئله وارد کنیم، می‌توانیم نتیجه بگیریم که دور منفی در این گراف وجود ندارد.) 

}


 % مسئلۀ ۱ب
%\زیرمسئله % مسئلۀ ۱پ
%\زیرمسئله % مسئلۀ ۱ت
\پایان{زیرمسائل}

%\solution{
%	پاسخ نمونه.
%}




%%%%%%%%%%%%%%%%%%%%%%%%%%%%%%%%%%%%%%%%%%%%%%%%%%%%%
% برخی دستورات مفید برای تمرین‌ها. %
%%%%%%%%%%%%%%%%%%%%%%%%%%%%%%%%%%%%%%%%%%%%%%%%%%%%%

% فرمول‌های نمونه:
%$f(x) = \Theta(x)$
%$T(x, y) \leq \log(x) + 2^y + \binom{2n}{n}$
% {\tt code\_function}


% لیست موارد بدون شماره‌گذاری:
%\شروع{فقرات}
%    \فقره مورد اول 
%        \شروع{فقرات}
%            \فقره مورد اول داخلی 
%                \شروع{فقرات}
%                    \فقره اولین مورد خیلی داخلی 
%                    \فقره دومین مورد خیلی داخلی 
%                \پایان{فقرات}
%            \فقره مورد دوم داخلی 
%        \پایان{فقرات}
%    \فقره مورد دوم 
%\پایان{فقرات}

% لیست موارد شماره‌گذاری شده:
%\شروع{شمارش}
%    \فقره یک مورد 
%    \فقره یک مورد دیگر 
%    \فقره باز هم یک مورد دیگر
%\پایان{شمارش}

% برای نوشتن یک فرمول چند خطی که یک شماره بخوره. برای حالت‌های دیگر به سایت زیر مراجعه کنید:
% https://www.overleaf.com/learn/latex/Aligning_equations_with_amsmath
%\begin{align} 
%    \begin{split}
%        (x+y)^3 &= (x+y)^2(x+y) \\
%                &= (x^2+2xy+y^2)(x+y) \\
%                &= (x^3+2x^2y+xy^2) + (x^2y+2xy^2+y^3) \\
%                &= x^3+3x^2y+3xy^2+y^3
%    \end{split}                                 
%\end{align}

% نوشتن ماتریس و آرایه دوبعدی
%\begin{align}
%    A = 
%    \begin{bmatrix}
%        A_{11} & A_{21} \\
%        A_{21} & A_{22}
%    \end{bmatrix}
%\end{align}

% اضافه کردن تصویر با شرح. برای اضافه کردن تصویر بدون شرح، فقط دستور درج‌تصویر را استفاده کنید.
%\شروع{شکل}[h!]
%\تنظیم‌ازوسط
%  \درج‌تصویر[پهنا=0.5\پهنای‌متن]{img.jpg}
%\شرح{شرح تصویر}
%\برچسب{تصویر:۱}
%\پایان{شکل}
%\پایان{زیرمسائل}

%کد زیر را ببینید.
%%\lstinputlisting{source_filename.py}
%\begin{latin}
%\begin{lstlisting}
%# Import the modules
%import sys
%import random
%
%ans = True
%
%while ans:
%question = raw_input("Ask the magic 8 ball a question: (press enter to quit) ")
%
%answers = random.randint(1,8)
%
%if question == "":
%sys.exit()
%
%elif answers == 1:
%print "It is certain"
%
%elif answers == 2:
%print "Outlook good"
%
%elif answers == 3:
%print "You may rely on it"
%
%elif answers == 4:
%print "Ask again later"
%
%elif answers == 5:
%print "Concentrate and ask again"
%
%elif answers == 6:
%print "Reply hazy, try again"
%
%elif answers == 7:
%print "My reply is no"
%
%elif answers == 8:
%print "My sources say no"
%\end{lstlisting}
%\end{latin}
